% =============================================================
% TRACEABILITY TABLE — Ch4 Roadmap Version
% Same core as Table 1.1 (tab:rq_obj_contrib_mapping) with a
% "Results Section" column instead of evidence.
% The companion table in Ch5 (traceability_table_v2.tex) adds
% the "Key Evidence" column after results have been presented.
% =============================================================
% Include via % =============================================================
% TRACEABILITY TABLE — Ch4 Roadmap Version
% Same core as Table 1.1 (tab:rq_obj_contrib_mapping) with a
% "Results Section" column instead of evidence.
% The companion table in Ch5 (traceability_table_v2.tex) adds
% the "Key Evidence" column after results have been presented.
% =============================================================
% Include via % =============================================================
% TRACEABILITY TABLE — Ch4 Roadmap Version
% Same core as Table 1.1 (tab:rq_obj_contrib_mapping) with a
% "Results Section" column instead of evidence.
% The companion table in Ch5 (traceability_table_v2.tex) adds
% the "Key Evidence" column after results have been presented.
% =============================================================
% Include via % =============================================================
% TRACEABILITY TABLE — Ch4 Roadmap Version
% Same core as Table 1.1 (tab:rq_obj_contrib_mapping) with a
% "Results Section" column instead of evidence.
% The companion table in Ch5 (traceability_table_v2.tex) adds
% the "Key Evidence" column after results have been presented.
% =============================================================
% Include via \input{traceability_table_ch4.tex}

\begin{landscape}
\begin{table}[htbp]
\centering
\caption{Research Roadmap: Mapping of Research Questions, Objectives, and Contributions to Results Sections}
\label{tab:rq_obj_results_roadmap}

\small
\setlength{\tabcolsep}{4pt}
\renewcommand{\arraystretch}{1.25}

\begin{tabularx}{\linewidth}{p{2.8cm} X X p{2.6cm} p{3.6cm}}
\toprule
\textbf{Research Question (RQ)} &
\textbf{Primary objective(s) that directly answer it} &
\textbf{Supporting objective(s) (enablers / infrastructure)} &
\textbf{Contribution(s) most directly produced} &
\textbf{Results section(s)} \\
\midrule

\textbf{RQ1:} How robust is MCQ-based evaluation to variations in distractor/option structure in systems engineering, and what risks does this pose when MCQ accuracy is used as primary qualification evidence? &
(5) Quantify MCQ robustness/sensitivity via statistical tests and effect sizes &
(1) Curate and version benchmark artifacts, including controlled MCQ variants suitable for robustness studies\newline
(3) Execute standardized inference across variants\newline
(7) Synthesize guidance framing robustness implications &
C2 (Robustness analysis of MCQ evaluation using distractor/option variation) &
\S\ref{sec:performance_mcq} MCQ Performance\newline
\S\ref{sec:comparison_evaluation_modalities} MCQ vs.\ OSQ Comparison \\
\midrule

\textbf{RQ2:} How do MCQ and OSQ modalities differ in the type and strength of evidence they provide for qualifying LLMs in systems engineering, and how can evidence from multiple modalities be combined to support task-aligned, risk-aware qualification decisions? &
(5) Quantify cross-modality evidence patterns and uncertainty\newline
(7) Synthesize qualification guidance (task-/risk-aligned fusion) &
(1) Curate and version benchmark artifacts, including paired MCQ/OSQ items suitable for cross-modality comparison\newline
(2) Generate OSQ grading scaffolds\newline
(3) Execute standardized inference\newline
(4) Score OSQ via rubric-based judging and judge stability checks &
C1 (Multi-format SysEngBench extension)\newline
C3 (Evaluation modality fusion as qualification basis) &
\S\ref{sec:conversion_success} Conversion Success\newline
\S\ref{sec:performance_osq} OSQ Performance\newline
\S\ref{sec:comparison_evaluation_modalities} MCQ vs.\ OSQ Comparison\newline
\S\ref{sec:modality_mapping} INCOSE Category \& Modality Mapping \\
\midrule

\textbf{RQ3:} How do evaluation modality and response length influence the computational cost of qualification, and what trade-offs arise between judged quality and token efficiency? &
(6) Evaluate token-efficiency and cost trade-offs; extend cost modeling to OSQ evaluation &
(3) Standardized inference with token/length logging\newline
(4) Rubric-based judging to produce quality signals &
C4 (Token-efficient qualification: cost-aware response length trade-offs) &
\S\ref{sec:tok} Tokenomics \& Cost \\
\bottomrule
\end{tabularx}

\vspace{4pt}
\noindent\textit{Note.} Chapter~\ref{ch:results} is organized by evaluation modality and analysis type rather than by objective. This table maps the objective-driven structure of Chapter~\ref{ch:method} to the results that follow. A companion table in Chapter~\ref{ch:discussion} (Table~\ref{tab:rq_obj_contrib_evidence}) revisits this mapping with the key evidence produced.
\end{table}
\end{landscape}


\begin{landscape}
\begin{table}[htbp]
\centering
\caption{Research Roadmap: Mapping of Research Questions, Objectives, and Contributions to Results Sections}
\label{tab:rq_obj_results_roadmap}

\small
\setlength{\tabcolsep}{4pt}
\renewcommand{\arraystretch}{1.25}

\begin{tabularx}{\linewidth}{p{2.8cm} X X p{2.6cm} p{3.6cm}}
\toprule
\textbf{Research Question (RQ)} &
\textbf{Primary objective(s) that directly answer it} &
\textbf{Supporting objective(s) (enablers / infrastructure)} &
\textbf{Contribution(s) most directly produced} &
\textbf{Results section(s)} \\
\midrule

\textbf{RQ1:} How robust is MCQ-based evaluation to variations in distractor/option structure in systems engineering, and what risks does this pose when MCQ accuracy is used as primary qualification evidence? &
(5) Quantify MCQ robustness/sensitivity via statistical tests and effect sizes &
(1) Curate and version benchmark artifacts, including controlled MCQ variants suitable for robustness studies\newline
(3) Execute standardized inference across variants\newline
(7) Synthesize guidance framing robustness implications &
C2 (Robustness analysis of MCQ evaluation using distractor/option variation) &
\S\ref{sec:performance_mcq} MCQ Performance\newline
\S\ref{sec:comparison_evaluation_modalities} MCQ vs.\ OSQ Comparison \\
\midrule

\textbf{RQ2:} How do MCQ and OSQ modalities differ in the type and strength of evidence they provide for qualifying LLMs in systems engineering, and how can evidence from multiple modalities be combined to support task-aligned, risk-aware qualification decisions? &
(5) Quantify cross-modality evidence patterns and uncertainty\newline
(7) Synthesize qualification guidance (task-/risk-aligned fusion) &
(1) Curate and version benchmark artifacts, including paired MCQ/OSQ items suitable for cross-modality comparison\newline
(2) Generate OSQ grading scaffolds\newline
(3) Execute standardized inference\newline
(4) Score OSQ via rubric-based judging and judge stability checks &
C1 (Multi-format SysEngBench extension)\newline
C3 (Evaluation modality fusion as qualification basis) &
\S\ref{sec:conversion_success} Conversion Success\newline
\S\ref{sec:performance_osq} OSQ Performance\newline
\S\ref{sec:comparison_evaluation_modalities} MCQ vs.\ OSQ Comparison\newline
\S\ref{sec:modality_mapping} INCOSE Category \& Modality Mapping \\
\midrule

\textbf{RQ3:} How do evaluation modality and response length influence the computational cost of qualification, and what trade-offs arise between judged quality and token efficiency? &
(6) Evaluate token-efficiency and cost trade-offs; extend cost modeling to OSQ evaluation &
(3) Standardized inference with token/length logging\newline
(4) Rubric-based judging to produce quality signals &
C4 (Token-efficient qualification: cost-aware response length trade-offs) &
\S\ref{sec:tok} Tokenomics \& Cost \\
\bottomrule
\end{tabularx}

\vspace{4pt}
\noindent\textit{Note.} Chapter~\ref{ch:results} is organized by evaluation modality and analysis type rather than by objective. This table maps the objective-driven structure of Chapter~\ref{ch:method} to the results that follow. A companion table in Chapter~\ref{ch:discussion} (Table~\ref{tab:rq_obj_contrib_evidence}) revisits this mapping with the key evidence produced.
\end{table}
\end{landscape}


\begin{landscape}
\begin{table}[htbp]
\centering
\caption{Research Roadmap: Mapping of Research Questions, Objectives, and Contributions to Results Sections}
\label{tab:rq_obj_results_roadmap}

\small
\setlength{\tabcolsep}{4pt}
\renewcommand{\arraystretch}{1.25}

\begin{tabularx}{\linewidth}{p{2.8cm} X X p{2.6cm} p{3.6cm}}
\toprule
\textbf{Research Question (RQ)} &
\textbf{Primary objective(s) that directly answer it} &
\textbf{Supporting objective(s) (enablers / infrastructure)} &
\textbf{Contribution(s) most directly produced} &
\textbf{Results section(s)} \\
\midrule

\textbf{RQ1:} How robust is MCQ-based evaluation to variations in distractor/option structure in systems engineering, and what risks does this pose when MCQ accuracy is used as primary qualification evidence? &
(5) Quantify MCQ robustness/sensitivity via statistical tests and effect sizes &
(1) Curate and version benchmark artifacts, including controlled MCQ variants suitable for robustness studies\newline
(3) Execute standardized inference across variants\newline
(7) Synthesize guidance framing robustness implications &
C2 (Robustness analysis of MCQ evaluation using distractor/option variation) &
\S\ref{sec:performance_mcq} MCQ Performance\newline
\S\ref{sec:comparison_evaluation_modalities} MCQ vs.\ OSQ Comparison \\
\midrule

\textbf{RQ2:} How do MCQ and OSQ modalities differ in the type and strength of evidence they provide for qualifying LLMs in systems engineering, and how can evidence from multiple modalities be combined to support task-aligned, risk-aware qualification decisions? &
(5) Quantify cross-modality evidence patterns and uncertainty\newline
(7) Synthesize qualification guidance (task-/risk-aligned fusion) &
(1) Curate and version benchmark artifacts, including paired MCQ/OSQ items suitable for cross-modality comparison\newline
(2) Generate OSQ grading scaffolds\newline
(3) Execute standardized inference\newline
(4) Score OSQ via rubric-based judging and judge stability checks &
C1 (Multi-format SysEngBench extension)\newline
C3 (Evaluation modality fusion as qualification basis) &
\S\ref{sec:conversion_success} Conversion Success\newline
\S\ref{sec:performance_osq} OSQ Performance\newline
\S\ref{sec:comparison_evaluation_modalities} MCQ vs.\ OSQ Comparison\newline
\S\ref{sec:modality_mapping} INCOSE Category \& Modality Mapping \\
\midrule

\textbf{RQ3:} How do evaluation modality and response length influence the computational cost of qualification, and what trade-offs arise between judged quality and token efficiency? &
(6) Evaluate token-efficiency and cost trade-offs; extend cost modeling to OSQ evaluation &
(3) Standardized inference with token/length logging\newline
(4) Rubric-based judging to produce quality signals &
C4 (Token-efficient qualification: cost-aware response length trade-offs) &
\S\ref{sec:tok} Tokenomics \& Cost \\
\bottomrule
\end{tabularx}

\vspace{4pt}
\noindent\textit{Note.} Chapter~\ref{ch:results} is organized by evaluation modality and analysis type rather than by objective. This table maps the objective-driven structure of Chapter~\ref{ch:method} to the results that follow. A companion table in Chapter~\ref{ch:discussion} (Table~\ref{tab:rq_obj_contrib_evidence}) revisits this mapping with the key evidence produced.
\end{table}
\end{landscape}


\begin{landscape}
\begin{table}[htbp]
\centering
\caption{Research Roadmap: Mapping of Research Questions, Objectives, and Contributions to Results Sections}
\label{tab:rq_obj_results_roadmap}

\small
\setlength{\tabcolsep}{4pt}
\renewcommand{\arraystretch}{1.25}

\begin{tabularx}{\linewidth}{p{2.8cm} X X p{2.6cm} p{3.6cm}}
\toprule
\textbf{Research Question (RQ)} &
\textbf{Primary objective(s) that directly answer it} &
\textbf{Supporting objective(s) (enablers / infrastructure)} &
\textbf{Contribution(s) most directly produced} &
\textbf{Results section(s)} \\
\midrule

\textbf{RQ1:} How robust is MCQ-based evaluation to variations in distractor/option structure in systems engineering, and what risks does this pose when MCQ accuracy is used as primary qualification evidence? &
(5) Quantify MCQ robustness/sensitivity via statistical tests and effect sizes &
(1) Curate and version benchmark artifacts, including controlled MCQ variants suitable for robustness studies\newline
(3) Execute standardized inference across variants\newline
(7) Synthesize guidance framing robustness implications &
C2 (Robustness analysis of MCQ evaluation using distractor/option variation) &
\S\ref{sec:performance_mcq} MCQ Performance\newline
\S\ref{sec:comparison_evaluation_modalities} MCQ vs.\ OSQ Comparison \\
\midrule

\textbf{RQ2:} How do MCQ and OSQ modalities differ in the type and strength of evidence they provide for qualifying LLMs in systems engineering, and how can evidence from multiple modalities be combined to support task-aligned, risk-aware qualification decisions? &
(5) Quantify cross-modality evidence patterns and uncertainty\newline
(7) Synthesize qualification guidance (task-/risk-aligned fusion) &
(1) Curate and version benchmark artifacts, including paired MCQ/OSQ items suitable for cross-modality comparison\newline
(2) Generate OSQ grading scaffolds\newline
(3) Execute standardized inference\newline
(4) Score OSQ via rubric-based judging and judge stability checks &
C1 (Multi-format SysEngBench extension)\newline
C3 (Evaluation modality fusion as qualification basis) &
\S\ref{sec:conversion_success} Conversion Success\newline
\S\ref{sec:performance_osq} OSQ Performance\newline
\S\ref{sec:comparison_evaluation_modalities} MCQ vs.\ OSQ Comparison\newline
\S\ref{sec:modality_mapping} INCOSE Category \& Modality Mapping \\
\midrule

\textbf{RQ3:} How do evaluation modality and response length influence the computational cost of qualification, and what trade-offs arise between judged quality and token efficiency? &
(6) Evaluate token-efficiency and cost trade-offs; extend cost modeling to OSQ evaluation &
(3) Standardized inference with token/length logging\newline
(4) Rubric-based judging to produce quality signals &
C4 (Token-efficient qualification: cost-aware response length trade-offs) &
\S\ref{sec:tok} Tokenomics \& Cost \\
\bottomrule
\end{tabularx}

\vspace{4pt}
\noindent\textit{Note.} Chapter~\ref{ch:results} is organized by evaluation modality and analysis type rather than by objective. This table maps the objective-driven structure of Chapter~\ref{ch:method} to the results that follow. A companion table in Chapter~\ref{ch:discussion} (Table~\ref{tab:rq_obj_contrib_evidence}) revisits this mapping with the key evidence produced.
\end{table}
\end{landscape}
